% Français 9115, batch 23 — Gallica views 441–460.
% Pass: Opus 5.5 (claude-opus-5-5), 2026-09-22 — first pass, unchecked against the views by a human.
%
% What the twenty views hold. One piece, begun before the batch and going
% on after it: the prize memoir on vibrating elastic surfaces, in French,
% in her regular fair hand, corrected more heavily as it goes (struck
% words, words above the line, notes of her own in the margins, the foot
% of several pages cancelled in superposed drafts). Pages (35) to (54) of
% her pagination, without a gap, one to a view. Leaves written on both
% sides, the ink number on the side her pagination makes odd, as in batch
% 22. No blank view, no microfilm target, no printed matter, no repeat of
% a photographed page. It is § 4, « Comparaison dela théorie à
% l'experience; dans le cas des plaques quarrées », begun at view 439: the
% nodal figures of the square plate with free edges, taken from equation
% (C), z = … {scriptA f(x) + B f(y)}, described value of ω by value of ω
% (first to fifth), and compared with Chladni's figures.
%
%   v. 441  (35) end of the sentence of v. 440: n² « intersections
%           fictives », hence n² invariable points of rest whatever the
%           distortion of the nodal figures of (C); back to B = −scriptA;
%           a general remark: what equal and symmetric portions of the
%           simple lamina between nodes would give on the diagonal.
%   v. 442  (36) under that hypothesis, n−1 (or n) nodal lines parallel
%           and perpendicular to the diagonals; not so in Euler's theory,
%           « entièrement confirmée par l'expérience », because z is
%           greatest at the ends of the simple lamina.
%   v. 443  (37) N.o 11, the cases one by one; « 1.ere valeur de l'angle
%           ω »: two nodes on the lamina, a second nodal line starting at
%           the edge, straight and perpendicular to the first diagonal.
%   v. 444  (38) the two diagonals are the whole figure and pass through
%           the n² = 4 invariable points; « 2.eme valeur de l'angle ω »
%           (four struck lines on n² = 9); f(x), f(y) for the two
%           functions of (C); she will need figures. The foot of the page,
%           on keeping Chladni's figure numbers, is struck in two
%           interwoven drafts.
%   v. 445  (39) a figure (square abcd, diagonal ad, dotted « lignes
%           fictives », a closed nodal curve, points i and e), with her
%           « N.ta »: she redrew it because she had left half an interval
%           between the first dotted line and the edge. Then the
%           conventions: dotted lines for the nodal lines of scriptA = 0,
%           B = 0; i their intersections; a the origin, ab the axis of x,
%           ac of y; each figure « par le chiffre que porte son analogue
%           dans l'ouvrage de M.r Chladni »; « (V pl. jointe a ce M. fig. 67
%           c) ».
%   v. 446  (40) the curve ie'e'i through two invariable points, nearly
%           straight where it crosses the diagonal. The page stops two
%           thirds down, without a full stop.
%   v. 447  (41) « Cette remarque est generale »; the curve again; the
%           diagonal and the closed curve pass through the n² = 9 points;
%           « 3.eme valeur de l'angle ω » (« V. la pl. … fig 72 a »).
%   v. 448  (42) a figure: the square with both diagonals, four dotted
%           lines each way, sixteen points i, a closed rounded-square
%           curve; then the description of the curve for the third value.
%   v. 449  (43) the curve reaches a new invariable point; « à cause de la
%           symetrie dela figure » only the second diagonal remains to be
%           drawn, from the point c; diagonals and closed curve pass
%           through the n² = 16 points; « 4.eme valeur de l'angle ω »
%           (« V. la pl. … fig. 78 »).
%   v. 450  (44) two figures at the head: a square with diagonal ad,
%           dotted lines, points i and a sinuous curve lettered e, e', e'',
%           e'''; beside it a larger square cancelled with long oblique
%           strokes. Then the curve for the fourth value, interval by
%           interval.
%   v. 451  (45) the curve goes on (the top rewritten line by line); a
%           general remark: three possible relations, <, =, >, between the
%           greatest values of f(y) and f(x) in an interval, which decide
%           how the nodal curve turns.
%   v. 452  (46) the three cases: concavity towards the dotted lines
%           parallel to y; four branches crossing towards the centre of
%           the interval (in the left margin); concavity towards the lines
%           parallel to x. « Je reviens à la description de la figure
%           nodale »; the last six lines of the page struck.
%   v. 453  (47) the first closed curve returns to its point e; a second
%           curve from a second point e, through the fourth and fifth
%           dotted lines.
%   v. 454  (48) the second curve « beaucoup plus resserrée » (4 points i
%           against 6); for an even number of dotted lines the figure is
%           four equal and similar quarters; the n² = 25 points; « 5.eme
%           valeur de l'angle ω » opens at the foot.
%   v. 455  (49) the curve for the fifth value, interval by interval
%           (« parallellement à l'axe des y », « des x »); four closed
%           curves and the two diagonals make the whole figure.
%   v. 456  (50) through the n² = 36 points. « Les cas suivans n'ayant pas
%           été soumis à l'experience, je ne m'y arreterai pas »: their
%           figures would follow from comparing f(x) and f(y), which are
%           « celles de z pour la lame dont les extrêmités sont libres au
%           sens d'Euler »; the conditions every later figure's closed
%           curve obeys, the third value excepted.
%   v. 457  (51) the end of that list, heavily struck; the closed curve
%           passes through 12 (doubtful) invariable points, 8 in the case
%           of the third value (left margin, fig. 72 a). N.o 12: other
%           figures of the same integral (C) for the same ω, by changing
%           the ratio of scriptA and B; only the invariable points are
%           common to them all; Chladni's « flexion » quoted (« V. N.o 109
%           de son traité d'acoustique »), with a note of hers, in the
%           interline and the margin, that his definition draws the axis
%           through the highest or lowest points where it should pass
%           through the mean ones.
%   v. 458  (52) one flexion for the first value, 1 + ½ for the second,
%           2, 2 + ½, 3 for the fifth, and why (a margin note on the half
%           flexion); « combien la théorie … est conforme à l'expérience »:
%           Chladni's fig. 64, fig. 65 (with « la table p. 150 de
%           Chladni »); the ratio of scriptA to B sets the height of each
%           sinuosity. The text stops a third of the way down.
%   v. 459  (53) Chladni's figs 67, 67 b, 67 c, 72 a, 72 b, 73 a, 73 b set
%           against her values of ω; a struck parenthesis on Chladni's
%           renumbering of two figures; « l'experience n'en a montré
%           jusqu'ici qu'une petite partie », the primitive figure seen in
%           one case only, the second value of ω (« V. fig 67 de
%           Chladni »). The foot of the page is a tangle of struck drafts.
%   v. 460  (54) Chladni's figs 78 and 85, for her fourth and fifth values:
%           his plain and dotted lines combine into figures like hers; « Au
%           premier coup d'œil » theory and experiment seem to disagree,
%           but any line satisfying the conditions at the limits is « une
%           veritable extrêmité analytique », so the parts it separates may
%           vibrate as if isolated, provided they give the same sound. The
%           page stops mid-sentence on « leurs mouvemens se ».
%
% Seams. View 441 continues the sentence that stops at the foot of view
% 440 (« si on represente par n le nombre de / ces lignes »). Inside the
% batch each page runs on to the next, except at v. 446/447 (v. 446 stops
% without a full stop two thirds down; v. 447 opens « Cette remarque est
% generale »), v. 452/453 (the foot of v. 452 is struck; v. 453 opens a new
% sentence), and v. 458/459 and 459/460 (new paragraphs). View 460 stops
% mid-sentence on « lorsque leurs mouvemens se »; view 461 (page (55), ink
% 226, batch 24) opens « rapporteroient à la même valeur de l'angle ω;
% mais elles pourroient présenter des figures differentes entr'elles »,
% which completes it.
%
% The numbers on the leaves, and why no \folio{} is written. (1) Her
% pagination, in parentheses at the head, centre: (35) at v. 441 to (54)
% at v. 460, one page per view, without a gap. The odd pages are struck
% with a stroke, the even ones not, as in batch 22. Her 4 is formed like
% a « b » in (44) and (48). (2) The ink series at the top right, in the
% larger rounder hand, on the odd-paginated side only: 216 (v. 441), 217
% (443), 218 (445), 219 (447), 220 (449), 221 (451), 222 (453), 223 (455),
% 224 (457), 225 (459) — ten leaves, one number each, without a gap; none
% on the even views, the backs of those leaves. It continues 215 at v. 439
% and goes on to 226 at v. 461. Penned, not pencilled, as everywhere in
% this volume: no \folio{} is written. (3) Her paragraph numbers: N.o 11
% (v. 443), N.o 12 (v. 457), continuing N.o 10 at v. 439; her ordinals of
% the values of ω as headings: « 1.ere » (v. 443), « 2.eme » (v. 444),
% « 3.eme » (v. 447), « 4.eme » (v. 449), « 5.eme » (v. 454); and a
% reference « § 3 », doubtful, to an earlier section (v. 460). (4)
% Chladni's figure numbers, which she adopts: 64, 65, 67, 67 b, 67 c, 72
% a, 72 b, 73 a, 73 b, 78, 85, and his N.o 109 and « table p. 150 ».
% Every number above was read on its view. None of these leaves can be
% Laubenbacher and Pengelley's « Manuscript C » (ff. 348r–349r): the
% subject is elasticity and the ink series runs 216–225.
%
% Notation. Hers, kept. scriptA is the looped cursive capital of views
% 439–440, distinct from B and from the side A of the square, set
% $\mathcal{A}$ as in batch 22. f(x), f(y) are her names for the two
% functions of (C); z, x, y, n, ω, i, e, a, c, ab, ac are underlined in
% the prose and keep \underline{} inside the maths; the underlined
% « flexion(s) » is set in \emph{}. Primes on the points e', e'', e''' are
% small strokes, read with some doubt (v. 446–450). Her « < que », « = »,
% « > que » are kept as signs. Ordinals written with a superscript
% (« 2.eme », « 1.ere », « N.ta ») are set with \textsuperscript. The
% figures are not drawn; each is described in a note with its lettering.
% Quotation of Chladni with « at the head of every line (v. 457). Spelling
% is hers and is not flagged: « rappeller », « rapelle », « parallelle »,
% « complette », « systême », « symetrie », « raproche », « resultans »,
% « suivans », « raisonnemens », « mouvemens », « tems », « longtems »,
% « fesant », « exposés » for exposées, « trouvé » for trouvées, « plus
% part », « aux quels », « cette intervalle », « isolés » and « isolées »
% on one page, and the run-together words « dela », « audelà »,
% « aureste », « afait », « aupoint », « cidessus », « demême »,
% « desorte », « deplus », « parceque », « parconséquent », « àprésent »,
% « jusqu'ici ».
%
% What could not be read, and what is offered with doubt. \ill{} stands
% eight times: a struck word after « j'ai » in the cancelled foot of
% v. 444; a small struck sign after « devra avoir » (v. 447); the group
% after « pl. » in « (V. la pl. … fig 72 a) » (v. 447); a struck word
% after « celui » and the end of the added word « tour… » (v. 452); on
% v. 457 the end of the fourth struck line and a blackened word after « un
% seul »; the end of the struck parenthesis on v. 459. \uncertain{} stands
% for « fictives » and « consécutives » (v. 442), « on » (v. 442),
% « appartenantes » and « suivi » (v. 444, struck), « refait » in the note
% of v. 445 and the letter after « fig. 67 » (given as « c », v. 445),
% « jointe au M. » (v. 449), « s'avancera » (v. 450), « sera » (v. 452),
% « que » (v. 454), « 12 » and « lorsqu'il s'agit » (v. 457), « tirer »
% (v. 457), « que j'ai exposée » (v. 458), « ce qui » and « faire croire »
% (v. 459, struck), « § 3 » (v. 460).
%
% Nothing here was checked against any published edition of Chladni, of
% Euler or of Sophie Germain's papers.
\documentclass[11pt,a4paper]{article}
% The preamble shared by every transcription of the mathematicians' manuscripts in Gallica.
%
% It rests on one idea: **uncertainty compounds, it does not resolve.** A
% transcription that smooths over a doubtful word has destroyed the only thing
% it was made for — a reader must be able to tell, at every point, what was on
% the page and what was supplied. The apparatus macros below are therefore the
% critical apparatus, not decoration.
%
% The same commands are recognised by `scripts/render.mjs`, which produces the
% site's reading view, and by `scripts/tei.mjs`. A macro added here must be
% added there too, or rendering fails loudly — which is the wanted behaviour.
%
% Comments here are English, like the rest of the repository. The documents
% this preamble sets are French, and that is deliberate: see the skills.

\ProvidesPackage{germain}[2026/09/15 Transcriptions of mathematicians' manuscripts in Gallica]

\usepackage{amsmath,amssymb,amsthm}
\usepackage{mathrsfs}
% Kept from the parent preamble so the renderer's diagram support stays
% available; these manuscripts draw no commutative diagrams, and nothing here uses it.
\usepackage{tikz-cd}
\usepackage[english,french]{babel}
\usepackage{fontspec}

% --- Greek ------------------------------------------------------------------
% The text face stays fontspec's default, Latin Modern, which has no Greek:
% XeTeX drops every glyph it lacks with a line in the log and nothing on the
% page, and the Greek words the manuscripts quote vanished from the PDFs. So
% the Greek and Greek Extended blocks get a XeTeX character class of their own,
% and entering or leaving a run of it switches the family to CMU Serif — the
% same Computer Modern design, with polytonic Greek — and back. Only the
% family changes: a Greek word in \textit or \textbf keeps its shape and series.
% The transitions to and from every other class carry the tokens that class
% has towards an ordinary letter, so French spacing before « ; » or inside
% « » still holds after a Greek word. No grouping, so no brace or \endgroup
% can come out unbalanced; maths is left alone, since XeTeX inserts nothing
% there. Kept identical in `exercices.sty`.
\makeatletter
\newfontfamily\germainGreekfont{cmun}[
  Extension=.otf, NFSSFamily=germaingreek,
  UprightFont=*rm, ItalicFont=*ti, SlantedFont=*sl,
  BoldFont=*bx, BoldItalicFont=*bi, BoldSlantedFont=*bl]
\newXeTeXintercharclass\germain@greek
\def\germain@greekrange#1#2{%
  \@tempcnta=#1\relax
  \loop
    \XeTeXcharclass\@tempcnta=\germain@greek
    \ifnum\@tempcnta<#2\advance\@tempcnta\@ne\repeat}
\germain@greekrange{"0370}{"03FF}
\germain@greekrange{"1F00}{"1FFF}
\def\germain@greekprev{\rmdefault}
\def\germain@greekin{%
  \edef\germain@greekprev{\f@family}\fontfamily{germaingreek}\selectfont}
\def\germain@greekout{\fontfamily{\germain@greekprev}\selectfont}
\def\germain@greekjoin{%
  \XeTeXinterchartokenstate=\@ne
  \@tempcnta=\z@
  \loop
    \ifnum\@tempcnta=\germain@greek\else
      \XeTeXinterchartoks\germain@greek\@tempcnta=\expandafter{%
        \expandafter\germain@greekout\the\XeTeXinterchartoks\z@\@tempcnta}%
      \XeTeXinterchartoks\@tempcnta\germain@greek=\expandafter{%
        \the\XeTeXinterchartoks\@tempcnta\z@\germain@greekin}%
    \fi
    \ifnum\@tempcnta<4095\advance\@tempcnta\@ne\repeat}
% babel-french sets its own classes, and saves and restores the interchar
% state, each time a language is selected: join again after it does.
\addto\extrasfrench{\germain@greekjoin}
\addto\extrasenglish{\germain@greekjoin}
\AtBeginDocument{\germain@greekjoin}
\makeatother

\usepackage{geometry}
\geometry{a4paper,margin=25mm,marginparwidth=28mm}
\usepackage{marginnote}
\usepackage{xcolor}
\usepackage[normalem]{ulem}
\setlength{\emergencystretch}{3em}
\usepackage{hyperref}
\hypersetup{colorlinks=true,linkcolor=black,urlcolor=black,pdfborder={0 0 0}}

% --- Batch metadata ---------------------------------------------------------
% Taken as they stand by the rendering script, which shows them at the head of
% the reading view. They fix what the file claims to cover. The macro names are
% the parent project's, so that the scripts stay shared; \folder here means the
% volume's slug — `fr-9115` — and \pages counts Gallica views.

\newcommand{\folder}[1]{\gdef\grothFolder{#1}}
\newcommand{\batch}[1]{\gdef\grothBatch{#1}}
\newcommand{\pages}[2]{\gdef\grothFirst{#1}\gdef\grothLast{#2}}
\newcommand{\foldertitle}[1]{\gdef\grothFoldertitle{#1}}

% The holder's own shelfmark, printed as they write it: « Français 9115 ».
\newcommand{\shelfmark}[1]{\gdef\grothShelfmark{#1}}

% The dating, copied verbatim from the holder's catalogue — never inferred
% here. « XIXe siècle » is the BnF's; a pass that dates a leaf from its content
% says so in a \note{}, as its own inference.
\newcommand{\dating}[1]{\gdef\grothDating{#1}}

% The status notice, printed in the title block and repeated as a diagonal
% watermark on every page of the PDF. The manuscripts are in the public
% domain; what this line declares is not a right but a quality — an
% unverified first pass by a machine. The skills write the line; a file
% without it gets no watermark, so leaving it out is visible in the output.
\newcommand{\watermark}[1]{\gdef\grothWatermark{#1}}

\gdef\grothFolder{?}\gdef\grothBatch{}\gdef\grothFirst{?}\gdef\grothLast{?}%
\gdef\grothFoldertitle{}\gdef\grothShelfmark{}\gdef\grothDating{}\gdef\grothWatermark{}

% --- The résumé -------------------------------------------------------------
\newenvironment{resume}
  {\small\setlength{\parskip}{0.45em}}
  {\par\medskip\hrule\medskip}

% --- Keywords ---------------------------------------------------------------
% In English, closing the modernised reading's résumé: the modern vocabulary
% under which to search for what the leaves construct. The manifest extracts
% them to tag the volume — this line is the single source of the tags.
\newcommand{\keywords}[1]{\par\smallskip\noindent
  {\footnotesize\textsc{Keywords} --- \textit{#1}}\par}

% --- The critical apparatus -------------------------------------------------

% The Gallica view number, in the margin. This is the anchor that lets a
% disagreement over a formula be settled by looking at *one* image rather than
% a volume of 750. The site uses it to turn the facsimile as the transcription
% is scrolled. A view is one scanned image: usually one side of a leaf.
\newcommand{\page}[1]{%
  \marginnote{\footnotesize\color{gray!70}\textsf{v.\,#1}}%
  \label{page:#1}\ignorespaces}

% The BnF's pencilled foliation, where the leaf carries it and a pass can read
% it: « 198r ». This is what the literature cites — « ff. 198r–208v » — and
% the one line that turns a citation into a view. Recorded, never inferred:
% a folio a pass cannot read is not written.
\newcommand{\folio}[1]{%
  \marginnote{\footnotesize\color{gray!50}\textsf{f.\,#1}}\ignorespaces}

% The modernised reading's coarser counterpart to \page{}: which views a
% section is drawn from.
\newcommand{\pagerange}[2]{%
  \marginnote{\footnotesize\color{gray!70}\textsf{v.\,#1--#2}}%
  \ignorespaces}

% Illegible: never guessed.
\newcommand{\ill}{\textcolor{red!60!black}{[\,\ldots\,]}}

% A doubtful reading: the word is given, and flagged as uncertain.
\newcommand{\uncertain}[1]{\uline{#1}}

% An editorial addition — a missing letter, an implied word.
\newcommand{\add}[1]{\textcolor{blue!50!black}{[#1]}}

% Struck out by the writer. What was crossed out is part of the document.
\newcommand{\struck}[1]{\sout{#1}}

% The transcriber's note, on the state of the page or the mathematics.
\newcommand{\note}[1]{\par\smallskip{\small\color{gray!60!black}\textit{[#1]}}\par\smallskip}

% A marginal note of hers, distinct from the body of the text.
\newcommand{\marginal}[1]{\par\smallskip\noindent\hspace{1em}%
  {\small\color{gray!60!black}#1}\par\smallskip}

% --- Title ------------------------------------------------------------------
% The title block is French, like the documents themselves.

\AddToHook{shipout/background}{%
  \ifx\grothWatermark\empty\else
    \begin{tikzpicture}[remember picture, overlay]
      \node[rotate=52, scale=3.4, align=center, text=gray!18,
            font=\sffamily\bfseries]
        at (current page.center) {\grothWatermark};
    \end{tikzpicture}%
  \fi}

\AtBeginDocument{%
  \begin{center}
    {\large\bfseries Manuscrits de mathématiciens (Gallica) ---
      \ifx\grothShelfmark\empty\grothFolder\else\grothShelfmark\fi}\\[2pt]
    {\itshape\grothFoldertitle}\\[4pt]
    {\small vues \grothFirst--\grothLast
      \ifx\grothBatch\empty\else\ (lot \grothBatch)\fi}\\[2pt]
    \ifx\grothDating\empty\else
      {\small\color{gray!60!black}Datation du catalogue : \grothDating}\\[2pt]
    \fi
    \ifx\grothWatermark\empty\else
      {\small\bfseries\color{red!45!gray}\grothWatermark}\\[2pt]
    \fi
    {\footnotesize\color{gray!60!black}Transcription établie d'après les images IIIF de Gallica.
     Source gallica.bnf.fr / Bibliothèque nationale de France.\\
     Les lectures incertaines sont soulignées, les illisibles notées [\,\ldots\,],
     les ajouts entre crochets ; \textsf{v.} numérote les vues de Gallica,
     \textsf{f.} la foliotation au crayon de la BnF.}
  \end{center}
  \vspace{1em}}

\endinput


\folder{fr-9115}
\shelfmark{Français 9115}
\batch{23}
\pages{441}{460}
\dating{1801-1900}
\watermark{Lecture automatique\\non vérifiée}
\foldertitle{Papiers de Sophie GERMAIN. — Recueil de dissertations et problèmes mathématiques et physiques.}

\begin{document}

\note{Numérotation des feuillets. Chaque page porte en tête, au milieu,
entre parenthèses, sa pagination : (35) à (54) aux vues 441 à 460, sans
lacune ; les pages impaires sont barrées d'un trait. Les feuillets sont
écrits des deux côtés, et seul le côté de pagination impaire porte en
haut à droite, à l'encre, d'une écriture plus ronde, un nombre de la
série qui court dans tout le volume : 216 (vue 441), 217 (443), 218
(445), 219 (447), 220 (449), 221 (451), 222 (453), 223 (455), 224 (457),
225 (459). Cette série est tracée à la plume et non au crayon, et aucun
« f. » n'est donc porté. Tous ces nombres ont été lus sur la vue ; aucun
n'est déduit.}

\page{441}

ces lignes, il y aura $\underline{n}^2$ intersections fictives; et on pourra
être sûr que, \struck{qu} \add{quelles} que soient les distorsions dont les figures
nodales renfermées dans l'intégrale (C) puissent être susceptibles,
il y aura toujours sur la surface $\underline{n}^2$ points de repos
invariables.

Je reviens à la supposition $B = -\mathcal{A}$; mais avant d'entrer
dans l'examen des differens cas qui appartiennent à chaque
valeur particulière de l'angle $\underline{\omega}$, je vais encore presenter une
remarque générale qui m'aidera beaucoup dans les recherches
suivantes:

Si on adoptoit pour un instant l'hypothese de l'égalité
d'étendue et de courbure entre les differentes portions dela
simple lame terminées par deux points de repos consécutifs
(hypothese dont on n'avoit pas encore remarqué l'inexactitude)
on trouveroit qu'en conservant à $\underline{x}$ une valeur correspondante
à un des points dela ligne nodale diagonale, il y auroit
dans l'espace séparé de ce point par $2, 4, 6\ldots$ lignes \struck{diag}
nodales fictives, deux valeurs de $y$ qui, parcequ'elles donneroient,
pour le cas $\mathcal{A}=0$, chacune une valeur de $\underline{z}$ égale
et de même signe que celle qui appartiendroit au point
pris sur la diagonale, devroient pour la présente valeur
de $\mathcal{A}$ rendre $\underline{z}$ nulle. Il est clair que pour la valeur
de $\underline{x}$ qui appartiendroit aupoint d'origine dela diagonale,
les deux valeurs de $y$ dont il s'agit se reduiroient à une
\note{En tête, au milieu, « (35) », barré d'un trait oblique ; en haut à
droite, à l'encre, d'une écriture plus ronde, « 216 ». La première ligne
achève la phrase de la vue 440 (« si on represente par $\underline{n}$ le
nombre de / ces lignes »). Les deux « $n$ » soulignés portent un petit 2
au-dessus : lu $\underline{n}^2$. Après « que, », deux lettres sont barrées
et « quelles » est écrit au-dessus. Le « (C) » de « l'intégrale » est
écrit en exposant, au-dessus de la fin du mot. $\mathcal{A}$ est la
capitale cursive bouclée des vues 439–440, distincte de $B$. Après
« lignes », un mot commencé, « diag », est barré. « appartiendroit » est
écrit sur une première forme du mot et soudé à « aupoint ». La phrase se
poursuit à la vue 442.}

\page{442}

seule qui conviendroit au point egalement éloigné de l'un et
l'autre \struck{points} \add{lignes} de repos \uncertain{fictives} \uncertain{consécutives}. Ainsi, dans le cas
ou il y auroit $2\underline{n}$ points de repos sur la simple lame,
la suite des points de repos qui résulteroient dela
supposition cidessus, formeroient $\underline{n}-1$ lignes nodales \struck{diagonales}
paralleles à la première diagonale et $\underline{n}-1$ lignes nodales
paralleles à la seconde; et, lorsqu'il y auroit $2\underline{n}+1$
points de repos sur la simple lame, on trouveroit
demême $\underline{n}$ lignes nodales paralleles à la diagonale et
$\underline{n}$ autres lignes nodales perpendiculaires aux premières,
et qui \add{les} couperoient \struck{les premières} aux limites dela
surface.

Les mêmes résultats ne peuvent subsister dans l'hypothese
déduite dela théorie d'Euler, théorie qui est entièrement
confirmée par l'expérience, et à laquelle l'application
aux cas que je crois pouvoir y ramener, ajoute encore
un nouveau degré d'évidence. D'abord il n'est plus vrai
en général que les systèmes des lignes \add{droites} paralleles et perpendiculaires
aux diagonales, \struck{ayant leur} ou plutôt les courbes nodales
qui les remplacent, ayant leurs origines \uncertain{aux} points de
limites dela surface: car il est aisé de prouver qu'alors
les valeurs de $\underline{z}$ aux \struck{l'} extrêmités dela simple lame,
sont plus grandes que dans aucun autre point de
la même lame; et \struck{il est aisé} \uncertain{on} voit que plus
\note{En tête, au milieu, « (36) », non barré ; aucun nombre en haut à
droite : c'est le dos du feuillet de la vue 441. La première ligne
achève la phrase de la vue 441 (« se reduiroient à une / seule »).
« points » est barré et « lignes » écrit au-dessus. Les fins de
« fictives » et de « consécutives » sont surchargées — une première
terminaison, apparemment masculine, accordée à « points », récrite
après la correction en « lignes » — et la lecture n'est pas assurée.
« diagonales » est barré après « nodales » ; « les » est ajouté
au-dessus de « couperoient », et « les premières » barré après ce mot ;
« droites » est ajouté au-dessus de « paralleles ». Dans « auroit »
(« lorsqu'il y auroit »), une première terminaison est surchargée.
Après « aux », un « l' » est barré. Dans la dernière ligne, « il est
aisé » est barré et le mot qui suit, écrit sur un premier mot, est lu
avec doute « on ». La phrase se poursuit à la vue 443.}

\page{443}

la différence est grande entre la valeur de $\underline{z}$ à l'origine
de la diagonale et la plus grande valeur de $\underline{z}$ dans
l'intervalle des deux \struck{nœuds} \add{lignes fictives de repos}, où se trouvoit, dans la première
hypothese, l'origine d'une ligne nodale parallelle à cette
diagonale et d'une autre ligne nodale perpendiculaire
à \struck{celle-ci} \struck{la parallelle} \add{même diagonale}; plus aussi l'origine de la courbe qui, dans
\struck{remplace} la vraie hypothese, remplace le systême de
ces deux droites, doit être eloignée du bord dela
surface.

N.o 11 Je viens enfin à l'examen detaillé des differens cas
\struck{auxquels} \add{aux quels} donnent lieu les \struck{diverses} valeurs successives de
l'angle $\underline{\omega}$. \add{1.\textsuperscript{ere} valeur de l'angle $\underline{\omega}$.} En commençant donc par \add{le cas} le plus simple,
je suppose qu'en fesant $B=0$, il n'y ait que deux
nœuds dans le mouvement de la simple lame. Il
est clair qu'alors, en partant du point de double extrêmité
origine dela diagonale due à la supposition $x=y$,
on trouve dans l'espace séparé de ce point par deux
lignes fictives, un autre point pour lequel la valeur de
$\underline{z}$, dans la lame, est lamême qu'à l'origine, et que
ce point est la seconde extrêmité de la lame. Il
doit donc y avoir, dans ce cas, une seconde ligne nodale
ayant son origine à la limite de la surface; deplus
cette ligne doit être droite et perpendiculaire à la
\note{Le feuillet, plus petit, est monté plus bas sur sa feuille. En
tête, au milieu, « (37) », barré d'un trait oblique ; en haut à droite,
à l'encre, « 217 ». La première ligne achève la phrase de la vue 442
(« on voit que plus / la différence est grande »). « nœuds » est écrit
avec l'œ ouvert ; il est barré, et « lignes fictives de repos » écrit
au-dessus. Devant « ; plus aussi », « celle-ci » est barré sur la ligne
et, au-dessus, « la parallelle » barré puis « même diagonale » écrit à
la suite ; l'article n'est pas récrit. « remplace » est barré en tête de
la ligne suivante. Le « N.o 11 » est dans la marge, en tête de
l'alinéa. Dans la marge de gauche, en face de « donnent », « auxquels »
est barré et « aux quels » écrit au-dessus. La fin de « lieu » est
surchargée ; « diverses » est barré. « 1.\textsuperscript{ere} valeur de
l'angle $\omega$. » est écrit au-dessus de la ligne, après « l'angle
$\omega$. », comme un titre ajouté ; « le cas » au-dessus de « le plus
simple ». La phrase se poursuit à la vue 444.}

\page{444}

première diagonale, c'est adire de même figure et demême
direction que dans la première hypothese: car \struck{quelque} \add{quelle que} soit celle des
valeurs de l'angle $\underline{\omega}$, qui en permette l'existence \add{dela seconde diagonale}, \struck{elle} \add{cette ligne} traversera,
dans ces differentes parties, des intervalles égaux à ceux que
traversent les portions correspondantes dela première diagonale
qu'elle rencontre à angle droit au centre dela plaque.

Ces deux diagonales nodales sont les seules lignes de repos
qui restent les mêmes en passant de la première hypothese
à celle qui rend compte des phénomènes. Dans le cas
présent, elles composent, à elles seules, la figure nodale; et
il est clair que chacune d'elles rencontre deux intersections
fictives. Cette figure passe donc, comme le veut la
théorie, par les $\underline{n}^2 = 4$ points invariables de repos.

2.\textsuperscript{eme} valeur de l'angle $\omega$. \struck{Il y a dans ce cas trois points de
repos dans la simple lame et parconséquent $\underline{n}^2 = 9$ points
invariables de repos dans les figures \add{fictives} nodales qui resultent
ici dela même de cette valeur de l'angle $\omega$.} Afin de simplifier l'examen
du cas présent et de ceux qui vont suivre je representerai
par $f(x)$ et $f(y)$ les deux fonctions séparées de $\underline{x}$ et $\underline{y}$
qui entrent dans la valeur de $\underline{z}$ donnée par l'équation
(C); et pour suivre la formation des lignes de repos
qui se multiplient à mesure que l'angle $\underline{\omega}$ augmente,
je serai forcé d'avoir recours à l'emploi des figures.

\struck{Je conserverai, pour chaque valeur de l'angle $\omega$. les N.o que M.r
Chladni a placés a coté des figures \uncertain{appartenantes}} \struck{et quoique j'ai
\ill{} dans ces fig. des lignes que M.r Chladni n'a pas
\uncertain{suivi}, je leur conserverai les mêmes chiffres par lesquels elles
sont numerotées dans} \add{son} \struck{ouvrage. de M.r Chladni.}
\note{En tête, au milieu, « (38) », non barré ; aucun nombre en haut à
droite : c'est le dos du feuillet de la vue 443. La première ligne
achève la phrase de la vue 443 (« perpendiculaire à la / première
diagonale »). « quelque » est barré et « quelle que » écrit au-dessus ;
le long trait qui traverse « celle des » est un trait de liaison, non
une rature. « dela seconde diagonale » est ajouté au-dessus de
« l'existence » ; « elle » est barré et « cette ligne » écrit au-dessus.
« 2.\textsuperscript{eme} » est écrit en exposant en tête de la ligne,
devant « valeur de l'angle $\omega$ », qui fait titre comme à la vue 443.
Les quatre lignes qui suivent sont barrées d'un trait horizontal chacune ;
« fictives » y est écrit au-dessus de « nodales », et « ici dela même »,
souligné, en tête de la dernière ligne barrée, sous « résultent ». Le pied de la page porte un passage
entièrement biffé de traits ondulés, en deux rédactions enchevêtrées : la
première commence dans la marge de gauche (« Je conserverai, pour
chaque / valeur de l'angle $\omega$. les / N.o que M.r Chladni a placés a
coté des figures ») ; la seconde court dans le corps de la page (« et
quoique … dans son ouvrage ») et se termine, au-dessous, par « de M.r
Chladni. » barré. L'ordre donné ici est celui de l'œil ; il n'est pas
assuré. Le mot qui suit « j'ai » est barré et ne se lit pas ; « dont »
est écrit au-dessus de « que » ; « appartenantes » et « suivi » sont lus
avec doute. « son », d'une encre plus pâle au bout de la ligne, n'est pas
barré ; « ouvrage. » commence par un trait en croix. Les vues 445 et 448
portent des figures de plaques carrées.}

\page{445}

\note{En haut de la page, au-dessus du texte, une figure : un carré, les
angles marqués $a$ en haut à gauche, $c$ en bas à gauche, $d$ en bas à
droite (l'angle en haut à droite ne porte pas de lettre lisible) ; la
diagonale $ad$ tracée d'un trait plein ; des lignes ponctuées, deux
verticales et deux horizontales, parallèles aux côtés, et d'autres en
tirets plus fins près des bords ; une grande courbe fermée, repassée
plusieurs fois, qui entoure le centre en s'allongeant le long de l'autre
diagonale et coupe les lignes ponctuées ; de petites lettres, $e$ sur la
diagonale près de $a$ et près de $d$, $e$, $o$, $t$ (lecture incertaine)
le long du haut de la courbe, $i$ près de l'angle en haut à droite.}

\marginal{N.\textsuperscript{ta} j'ai \uncertain{refait} cette
fig parceque j'avois
laissé entre la première
ligne ponctuée et le
bord dela surface la
moitié de l'intervalle
entre deux lignes fictives
cequi ne doit pas être.}

J'indiquerai toujours par des lignes ponctuées les lignes \struck{n}nodales
qui se formeroient sur la plaque quarrée dans les suppositions
successives \struck{$f(x)=0$}, \struck{$f(y)=0$} \add{$\mathcal{A}=0$, $B=0$}; par $\underline{i}$, les points de repos
invariables résultans des intersections de ces lignes. \struck{Et} Je
\struck{supposerai tracée} prendrai \add{le point} $\underline{a}$ pour l'origine, \add{la ligne} $\underline{ab}$ pour l'axe
des $\underline{x}$, \add{la ligne} $\underline{ac}$ pour celui des $y$. Et je supposerai tracée la
diagonale qui indépendamment des valeurs particulières de l'angle
$\underline{\omega}$ resulte toujours de la supposition $y = x$ \add{lorsque $\mathcal{A} = -B$}; et je désignerai
chaque figure par le chiffre que porte son analogue dans l'ouvrage de M.r Chladni.
Cela posé, pour décrire, dans le cas présent, le reste de
la figure nodale; \add{(V pl. jointe a ce M. fig. 67 \uncertain{c})}, à partir de l'origine je passe audelà de
la seconde ligne nodale \add{fictive} parallelle à $\underline{ab}$; et je cherche
dans l'intervalle compris entre la seconde et la troisième lignes
\struck{nodales} \add{fictives, à quel point répond} \struck{quelle est} la plus grande valeur de $f(y)$. ce point $\underline{e}$
sera l'origine d'une ligne de repos, pour lequel les valeurs de $f(x)$ et
$f(y)$ seront les mêmes que pour le point $\underline{e}$ dela diagonale;
et il est clair que le \add{même} point $\underline{e}$ sera d'autant plus près dela
\note{En tête, au milieu, « (39) », barré ; en haut à droite, à l'encre,
« 218 ». La première ligne du texte ne suit pas la dernière de la vue
444, qui s'achève sur une phrase complète (« l'emploi des figures. ») et
un passage biffé ; la page s'ouvre sur la figure. La note marginale est
écrite à droite de la figure, ouverte par un signe en forme de dièse
suivi de « ta » en exposant, lu « N.\textsuperscript{ta} » comme à la
vue 334 ; le verbe, « refait », n'est pas assuré. Le « n » qui précède
« nodales » est barré et le mot repassé. Au-dessus de « $f(x)=0$,
$f(y)=0$ », barrés, est écrit « $\mathcal{A}=0$, $B=0$ » ; $\mathcal{A}$
est la capitale cursive bouclée des vues 439–441. « Et » est barré en
fin de ligne, puis « supposerai tracée » en tête de la suivante ; « le
point » est écrit au-dessus de « $a$ », « la ligne » au-dessus de
« $ab$ » et de « $ac$ ». « lorsque $\mathcal{A}=-B$ » est ajouté
au-dessus de « $y = x$ ». La ligne « chaque figure … de M.r Chladni »
est serrée entre deux lignes, d'une écriture plus petite, comme une
addition. Le renvoi « (V pl. jointe a ce M. fig. 67 c) » est écrit
au-dessus de « nodale ; » ; « ce M. » est surchargé d'une tache, et la
lettre après « 67 » n'est pas assurée. « fictive » est ajouté au-dessus
de « parallelle ». Devant « la plus grande », « nodales » et « quelle
est » sont barrés, et « fictives, à quel point répond » écrit au-dessus,
avec une tache sur la fin de « répond ». « même » est ajouté au-dessus de
« point ». Les deux points sont désignés par la même lettre $e$,
soulignée. La phrase se poursuit à la vue 446.}

\page{446}

limite de la surface, que \struck{$f(x$} la valeur de $f(y)$ sera plus
grande; \struck{pour les points suivans $e'$ dela diagonale} on
trouvera dans l'intervalle des seconde et troisième lignes
fictives, et de part et d'autre du point $e$, des points
$e'$ pour lesquels $f(y)$ et $f(x)$ auront les mêmes valeurs
que pour le point $e'$ de la diagonale; et la suite des
points qui auront la même relation avec les points
suivans de la diagonale, formera une ligne courbe
$ie'e'i$ qui passera par deux des points invariables de
repos. \struck{Audelà} \add{En deça} d'un de ces points et \struck{en deça} dela
seconde ligne fictive, la courbe nodale se trouvera
dans un seul et même intervalle avec la ligne nodale
diagonale qu'elle traversera sous une ligne presque droite,
parceque les valeurs successives de $f(y)$ prises pour les
points de la diagonale et pour les points correspondans
dela ligne nodale, ne peuvent \struck{que} differer beaucoup
entr'elles \struck{à cause qu'elles appartiennent} \add{appartenant} l'une et l'autre
à une seule et même branche de la courbe à laquelle
se plieroit la simple lame; desorte même que ces valeurs
seroient égales, si la courbure des deux portions d'une
partie de la simple lame, comprise entre deux nœuds,
étoit symétrique par rapport au point \struck{pour lequel}
qui répond à la plus grande valeur de $f(y)$
\note{En tête, au milieu, « (40) », non barré ; aucun nombre en haut à
droite : c'est le dos du feuillet de la vue 445. La première ligne
achève la phrase de la vue 445 (« d'autant plus près dela / limite de la
surface »). Après « que », un « $f(x$ » commencé est barré. « pour les
points suivans $e'$ dela diagonale » est barré d'un trait ondulé. Les
lettres des points sont petites et accentuées d'un trait : lues $e'$ ;
la courbe est désignée par quatre lettres, lues « $ie'e'i$ », sans
assurance sur les accents. « Audelà » est barré et « En deça » écrit
au-dessus ; plus loin « en deça » est barré. « que » est barré devant
« differer » ; « à cause qu'elles appartiennent » est barré et
« appartenant » écrit au-dessus. « nœuds » est écrit avec l'œ ouvert.
« pour lequel » est barré en fin de ligne. La page s'arrête aux deux
tiers, sur « $f(y)$ », sans point ; la vue 447 s'ouvre sur une phrase
nouvelle, « Cette remarque est generale ».}

\page{447}

Cette remarque est generale: ainsi je me contenterai dela rappeller
dans les cas semblables qui se rencontreront en examinant les
figures relatives aux valeurs suivantes de l'angle $\underline{\omega}$.

Je reviens à la description dela courbe nodale; après
avoir coupé la première ligne fictive parallelle à $\underline{ab}$
dans un des points $\underline{i}$, cette courbe s'éloignera dela ligne
fictive, en se rapprochant dela limite \add{dela surface} jusqu'au point $\underline{e''}$
pour lequel $f(x)$ aura la plus grande valeur que cette
même fonction puisse atteindre dans l'intervalle entre
la seconde et la troisième lignes fictives, \struck{pour} \add{car} ce point
$e''$ \struck{devra avoir} \struck{\ill{}} $f(x)$ et $f(y)$ devront avoir les mêmes
valeurs que pour le point correspondant dela diagonale.

Il est inutile d'achever la description dela figure; \struck{car} \add{puisqu'} à cause
de la symetrie de ses parties, on retomberoit dans les
\struck{mêmes} raisonnemens que j'ai exposés plus haut. On voit aureste
que l'ensemble de la ligne nodale diagonale et de
la \struck{courbe} courbe rentrante qui forment la figure
nodale, passe, comme le veut la théorie par les
$\underline{n}^2 = 9$ points \struck{de repos} invariables de repos

3.\textsuperscript{eme} valeur de l'angle $\omega$. Pour décrire la figure qui
convient a ce cas, \add{(V. la pl. \ill{} fig 72 a)}, à partir de l'origine je passe
audelà dela seconde ligne fictive parallelle à $\underline{ab}$,
\note{Le feuillet, plus petit, est monté plus bas sur sa feuille. En
tête, au milieu, « (41) », barré ; en haut à droite, à l'encre, « 219 ».
La page ne continue pas la phrase de la vue 446, qui s'arrête sans point
sur « la plus grande valeur de $f(y)$ » : elle s'ouvre sur une phrase
nouvelle, « Cette remarque est generale ». « dela surface » est ajouté
au-dessus de « limite » ; « pour » est barré et « car » écrit au-dessus.
Après « $e''$ », « devra avoir » est barré, puis un petit signe en croix
qui n'est pas lu. Devant « à cause », « car » est barré et « puisqu' »
écrit au-dessus ; « mêmes » est barré en tête de la ligne suivante, et
le premier « courbe », surchargé, l'est aussi. Après « points », « de
repos » est barré. La ligne des $\underline{n}^2 = 9$ points n'a pas de
point final. « 3.\textsuperscript{eme} » est écrit en exposant devant
« valeur de l'angle $\omega$ », qui fait titre comme aux vues 443 et 444.
Le renvoi entre parenthèses est écrit au-dessus de « a ce cas » ; après
« pl. », un groupe surchargé, peut-être un numéro de planche, ne se lit
pas. La phrase se poursuit à la vue 448.}

\page{448}

\note{En haut de la page, une figure : un carré aux angles marqués $a$ en
haut à gauche, $c$ en haut à droite, $b$ en bas à gauche, $d$ en bas à
droite ; la diagonale $ad$ d'un trait fort et la diagonale $bc$ d'un
trait fin ; quatre lignes verticales et quatre horizontales en tirets,
parallèles aux côtés, dont les croisements portent seize petits $i$ ;
une courbe fermée d'un trait fort, en forme de carré aux angles
arrondis posé sur la pointe, qui passe par les $i$ voisins du milieu des
côtés ; un $e$ sur la diagonale près de $a$ et un $e$ au sommet de la
courbe, en haut. C'est, d'après le texte, la figure de la troisième
valeur de l'angle $\omega$.}

jusqu'au point $\underline{e}$ qui donne la plus grande valeur
de $f(y)$ comprise entre la seconde et la troisième lignes
fictives; par les raisons exposées dans le cas precedent, ce
point sera l'origine d'une ligne de repos; et de part et
d'autre de cette origine se trouveront des suites de points
de repos qui formeront les deux branches d'une courbe
qui joindra deux points invariables de repos. Demême
encore audelà de celui de ces points qui se raproche
de l'origine, la ligne nodale se trouvera comprise
dans un seul et même intervalle avec la première
ligne nodale diagonale qu'elle coupera, en formant
dans cet intervalle une ligne presque droite. La courbe
nodale se raprochant toujours de la ligne $\underline{ab}$, elle
passera par un troisième point de repos invariable.
audelà de ce point elle n'atteindra pas tout afait
la plus grande valeur de $f(y)$; mais elle arrivera à
son sommet, et changera par conséquent de direction
\note{En tête, à gauche au-dessus de la figure, « (42) », non barré ;
aucun nombre en haut à droite : c'est le dos du feuillet de la vue 447.
La première ligne du texte achève la phrase de la vue 447 (« je passe
audelà dela seconde ligne fictive parallelle à $ab$, / jusqu'au point
$e$ »). Le « D » de « Demême » est repassé. La page se termine sur
« direction », suivi d'un trait qui peut être un point ; la phrase se
poursuit à la vue 449 (« de direction / lorsqu'elle sera parvenue »).}

\page{449}

lorsqu'elle sera parvenue a une valeur de $f(y)$ égale à la plus grande valeur
que la même fonction puisse recevoir entre la seconde et la troisième
lignes fictives; \struck{et} à partir de ce point, elle ira joindre diagonalement
et par une ligne peu differente de la droite, un nouveau point
invariable de repos. Mais à cause de la symetrie dela figure,
il est inutile de suivre plus loin sa description; il suffira
d'indiquer la formation de la seconde ligne nodale diagonale:
à partir de l'origine $\underline{a}$, je passe audelà de la quatrième ligne
fictive parallelle à $\underline{ab}$ jusqu'au point $\underline{c}$ qui donne une valeur
de $f(y)$ égale à celle dela même fonction au point $\underline{a}$. Ce point
sera donc l'origine d'une ligne nodale; et parceque pour chaque
point de la même ligne les valeurs de $f(y)$ seront toujours
égales aux valeurs correspondantes de la même fonction
pour la première diagonale, cette ligne sera necessairement
la seconde diagonale dont j'ai déja prouvé l'existence
au moyen de l'analyse.

En joignant le systême des deux diagonales à la courbe
nodale rentrante dont j'ai indiqué la formation, on voit
que la figure nodale passe, comme le veut la théorie,
par les $\underline{n}^2 = 16$ points invariables de repos.

4.\textsuperscript{eme} valeur de l'angle $\omega$. Pour décrire la figure qui convient
a ce cas, \add{(V. la pl. \uncertain{jointe au M.} fig. 78)} à partir de l'origine $\underline{a}$ je passe audelà de
la seconde ligne fictive parallelle à $\underline{ab}$ jusqu'au point
$\underline{e}$ auquel appartient la plus grande valeur de $f(y)$ comprise
entre la seconde et la troisième ligne \struck{ponctuée} \add{fictive}. Ce point
\note{En tête, au milieu, « (43) », barré ; en haut à droite, à l'encre,
« 220 ». La première ligne achève la phrase de la vue 448 (« et changera
par conséquent de direction / lorsqu'elle sera parvenue ») ; le « l »
initial est repassé d'un trait épais. Après « fictives; », un « et »
est barré sous l'« à » repassé ; l'« A » de « à partir » est aussi
repassé à la ligne « à partir de l'origine $a$ ». « 4.\textsuperscript{eme} »
est écrit en exposant devant « valeur de l'angle $\omega$ », qui fait
titre. Le renvoi entre parenthèses est écrit au-dessus de « a ce cas »,
d'une écriture très fine : après « pl. », un mot et une abréviation à
lettre haute sont lus avec doute « jointe au M. », comme à la vue 445 ;
le dernier chiffre de « 78 » est surchargé. « ponctuée » est barré et
« fictive » écrit au-dessus. La phrase se poursuit à la vue 450.}

\page{450}

\note{En haut de la page, deux figures côte à côte. À gauche, un carré
aux angles marqués $a$ en haut à gauche, $c$ en haut à droite, $b$ en
bas à gauche, $d$ en bas à droite ; la diagonale $ad$ ; des lignes
horizontales et verticales en tirets, aux croisements desquelles sont
de petits $i$ ; plusieurs courbes sinueuses, et le long d'elles des
lettres lues $e$, $e'$, $e''$, $e'''$ ; une tache d'encre près de
l'angle $d$. À droite, un carré plus grand, aux lignes en tirets, avec
une diagonale, une courbe sur la gauche et un trait vertical plein près
du côté droit, entièrement couvert de longs traits obliques
parallèles : une figure annulée.}

sera le sommet d'une ligne courbe nodale qui joindra deux
points de repos invariables. \struck{audela} En deça de celui de ces points
qui appartient à la seconde ligne fictive, la courbe nodale
se trouvera dans un seul et même intervalle avec la diagonale
qu'elle coupera parconséquent, en formant, \struck{une ligne peu differente}
\struck{de la droite} dans cette portion de son cours, une ligne peu
differente d'une droite. \struck{Elle en sortant de cet interva en deça} \add{A la limite} de
cet intervalle, elle coupera la première ligne fictive \struck{et} \add{en un point $\underline{i}$,} s'éloignera
ensuite de cette ligne en se rapprochant du bord dela plaque
jusqu'au point pour lequel $f(x)$ aura la plus grande valeur
que puisse atteindre la même fonction entre la seconde et
la troisième lignes fictives. Audelà de ce point, la courbe nodale
changera de direction et se rapprochera de la première
ligne fictive qu'elle coupera une seconde fois dans un
point $\underline{i}$. A partir de ce point, la courbe nodale \uncertain{s'avancera}
dans l'intervalle entre la première et la seconde lignes
fictives parallelles à $\underline{ab}$ jusqu'à ce qu'elle ait atteint
\note{En tête, au milieu, « (44) », non barré, le 4 formé comme un « b » ;
aucun nombre en haut à droite : c'est le dos du feuillet de la vue 449.
La première ligne achève la phrase de la vue 449 (« Ce point / sera le
sommet »). Le premier mot barré après « invariables. » est lu
« audela ». Après « en formant, », « une ligne peu differente de la
droite » est barré sur deux lignes. « Elle en sortant de cet interva »
est barré, « en deça » aussi, et « A la limite » est écrit au-dessus,
l'« A » repassé sur le « en » ; « de cet intervalle » reste. « et » est
barré devant « s'éloignera » et « en un point $i$, » écrit au-dessus.
Le verbe « s'avancera » est surchargé et n'est pas assuré. Le texte
occupe les deux tiers de la page sous les figures ; la phrase se
poursuit à la vue 451.}

\page{451}

la plus grande valeur que $f(y)$ puisse recevoir entre la troisième
et la quatrième lignes fictives; ou, cequi est la même chose,
entre la seconde et la troisième des mêmes lignes. Elle tournera
sa convexité vers la \struck{quatrième} ligne fictive parallelle à $\underline{ac}$; et à partir du point
\struck{rapprochera ensuite de la première ligne fictive qu'elle coupera}
auquel appartient la plus grande valeur de $f(y)$, elle se rapprochera de
\struck{encore dans un nouveau point de repos invariable de repos.}
la troisième qu'elle coupera encore dans un second point invariable de repos.

Afin de saisir plus aisément la \struck{formation} cause des \struck{sinuosités} \add{distorsions}
des lignes nodales, on peut observer qu'en
général il y a trois relations possibles entre la valeur
de $f(y)$ dans l'intervalle où l'on considère la courbe
nodale, et \struck{celle} \add{valeur} de la même fonction dans l'intervalle
où se trouve la partie correspondante dela diagonale
nodale. En effet on peut supposer que la première
valeur de $f(y)$ \struck{soit} \add{est} $<$, $=$, ou $>$ que la seconde.
\add{ensuite que,} \struck{et comme} pour les points dela diagonale $f(y)=f(x)$;
\add{et que pour les differens} \struck{et les} points de la courbe nodale, \struck{ont} les valeurs de $f(x)$
sont les mêmes que pour ceux qui leur correspondent dans
la diagonale. \struck{desorte} \add{On voit donc} que les relations dont il s'agit
existent \struck{exactement} \add{également} entre les valeurs de $f(y)$ et celles de $f(x)$
qui appartiennent à la portion que l'on considère dans la
courbe nodale. Cela posé, on voit aisément que si la
plus grande valeur de $f(y)$ dans l'intervalle où l'on examine
le cours de la courbe nodale est $<$ \add{que} la plus grande valeur
de $f(x)$ dans le même intervalle, la \struck{ligne} \add{courbe} nodale
\note{En tête, au milieu, « (45) », barré ; en haut à droite, à l'encre,
« 221 ». La première ligne achève la phrase de la vue 450 (« jusqu'à ce
qu'elle ait atteint / la plus grande valeur »). Le haut de la page est
une récriture ligne par ligne : deux lignes barrées d'un trait
horizontal (« rapprochera ensuite … qu'elle coupera » et « encore dans
un nouveau point … de repos. ») alternent avec les lignes gardées, qui
se lisent à la suite ; devant « ligne fictive parallelle à $ac$ », un
mot surchargé et barré est lu « quatrième ». « formation » est barré ;
« sinuosités » barré et « distorsions » écrit au-dessus ; « celle »
barré et « valeur » au-dessus ; « soit » barré et « est » au-dessus.
« ensuite que, » est écrit dans la marge de gauche devant « et comme »,
barré ; « et que pour les differens », dans la marge aussi, remplace
« et les », barré, au-dessus duquel est encore écrit « pour » ; « ont »
est barré devant « les valeurs ». « desorte » est barré et « On voit
donc » écrit au-dessus ; le mot barré devant « entre », lu
« exactement », est remplacé par « également ». Le « si » de « si la »
est repassé ; « que » est ajouté au-dessus de « $<$ la » ; « ligne » est
barré et « courbe » écrit au-dessus. Dans le dernier $f(x)$, le $x$ est
écrit sur un $y$. La phrase se poursuit à la vue 452.}

\page{452}

changera de direction, avant d'avoir atteint le point auquel
appartient la plus grande valeur de $f(x)$, qui \struck{est le}
\struck{point} point qui se trouve \struck{au centre} \add{vers le milieu} de l'intervalle forme
limité par \struck{les quatre lignes} \add{les deux lignes fictives consécutives parallelles a l'axe des $y$} fictives, \struck{consécutives chacune a sa}
\struck{parallelle}: Cette ligne tournera donc sa concavité vers une
\struck{des} de ces lignes parallelles à l'axe des $y$. Si la plus
grande valeur de $f(y)$ est égale à la plus grande valeur
de $f(x)$ dans le même intervalle, la courbe nodale \uncertain{sera}
\struck{se dirigera diagonalement vers le point $\underline{i}$ opposé a celui \ill{} par}
\struck{formée du systême de quatre branches qui se couperont}
\struck{lequel elle a passé en entrant dans l'intervalle.}
\struck{au centre de l'intervalle}
\marginal{formée du systême
de quatre branches
qui se couperont
vers le centre
de l'intervalle.}
Enfin, si la plus grande
valeur de $f(x)$ est $>$ que celle de $f(y)$, la courbe
nodale tournera sa concavité vers une des lignes
fictives parallelle à l'axe des $x$.

Je reviens à la description de la figure nodale.
Après avoir coupé la \struck{première} \add{troisième} ligne fictive \add{parallelle à $\underline{ac}$} dans \struck{le} un
\struck{quatrième} des points $\underline{i}$, \struck{qui lui appartiennent} la
courbe nodale \struck{s'approchera du bord dela surface} \add{tour\ill{} se trouvera dans un même intervalle avec}
la diagonale qu'elle traversera par une ligne presque droite,
\struck{jusqu'à ce qu'elle atteigne la valeur de $f(y)$ égale à la}
\struck{plus grande valeur de la même fonction dans l'intervalle}
\struck{entre la quatrième et la cinquième, ou cequi revient}
\struck{au même entre la première et la seconde ligne}
\struck{fictive. à partir de ce point elle se rapprochera encore}
\struck{dela première ligne fictive qu'elle coupera dans}
\note{Le feuillet, plus petit, est monté bas sur sa feuille. En tête, au
milieu, « (46) », non barré ; aucun nombre en haut à droite : c'est le
dos du feuillet de la vue 451. La première ligne achève la phrase de la
vue 451 (« la courbe nodale / changera de direction »). « atteint » est
traversé d'un trait qui peut être de liaison. Après « qui », « est le »
est barré, puis « point » en tête de la ligne suivante ; « au centre »
est barré et « vers le milieu » écrit au-dessus ; « forme » est laissé
en fin de ligne devant « limité ». « les quatre lignes » est barré et
« les deux lignes fictives consécutives parallelles a l'axe des $y$ »
écrit au-dessus, sur deux lignes ; « consécutives chacune a sa
parallelle » est barré. Après « la courbe nodale », un mot commençant
par « se » est surchargé et noirci ; « sera » est offert avec doute. Les
quatre lignes qui suivent sont barrées une à une ; elles mêlent deux
rédactions, et un mot surchargé après « celui » ne se lit pas. La fin
de la phrase est dans la marge de gauche (« formée du systême … de
l'intervalle. »). Plus bas, « première » est barré et « troisième »
écrit au-dessus ; « parallelle à $ac$ » est ajouté au-dessus de
« fictive » ; « le » est barré et « un » écrit après ; « quatrième » et
« qui lui appartiennent » sont barrés. « s'approchera du bord dela
surface » est barré, et au-dessus sont écrits un mot court, lu
« tour » suivi d'une lettre non lue, puis « se trouvera dans un même
intervalle avec ». Les six dernières lignes de la page sont barrées
d'un trait fin continu ; la vue 453 commence une phrase nouvelle.}

\page{453}

En sortant de cet intervalle, la courbe nodale passera par un point $\underline{i}$ et
se trouvera dans l'intervalle entre la première \add{et la seconde} lignes fictives parallelles
à $\underline{ac}$. Dans cet intervalle, la plus grande valeur de $f(y)$ sera
$>$ que la plus grande valeur de $f(x)$. ainsi, d'après ceque j'ai
dit plus haut, il est clair que la courbe nodale changera de
direction et se \struck{dirigera} \add{prolongera} parallellement à l'axe des $\underline{x}$. Elle
sortira de l'intervalle où je la considère àprésent pour rentrer
dans celui où se trouve le point $\underline{e}$ de départ. On voit donc
que cette partie dela figure nodale sera une courbe rentrante.
A partir du point $\underline{a}$, je passerai audelà dela quatrième ligne
fictive parallelle à $\underline{ab}$ jusqu'à un second point $\underline{e}$ qui sera le
sommet d'une nouvelle courbe nodale, et auquel appartiendra
parconséquent la plus grande valeur de $f(y)$ entre la première et la
seconde lignes fictives parallelles à $\underline{ab}$. La courbe, après avoir passé
par un point $\underline{i}$, \struck{passera} \add{se trouvera} dans l'intervalle entre la troisième et la
quatrième lignes fictives; et la portion renfermée dans cet intervalle,
sera semblable a celle qui, dans le même intervalle, appartient
à la première courbe nodale. En sortant de cet intervalle, la courbe
\struck{se} passera dans celui \struck{entre} \add{qui est formé par} la quatrième et la cinquième lignes fictives,
parceque cette courbe doit toujours être separée par un nombre
pair de lignes fictives, des points dela diagonale auxquels elle
correspond; et elle sera dirigée parallellement à l'axe des $y$, parceque
dans cet intervalle, la plus grande valeur de $f(y)$ est $<$ que
la plus grande valeur de $f(x)$. Après avoir passé par un
\note{Le feuillet est monté sur sa feuille. En tête, au milieu, « (47) »,
barré ; en haut à droite, à l'encre, « 222 ». La page s'ouvre sur une
phrase nouvelle, les dernières lignes de la vue 452 étant barrées. « et
la seconde » est ajouté au-dessus de « première » ; « dirigera » est
barré et « prolongera » écrit au-dessus ; « passera » est barré et « se
trouvera » écrit au-dessus ; « se » est barré devant « passera » ;
« entre » est barré et « qui est formé par » écrit au-dessus. Dans
« Dans cet intervalle », le « cet » est surchargé. « parallelles à
$ab$ », à la fin de la phrase sur le second point $e$, est souligné.
Les signes « $>$ que » et « $<$ que » sont dans la phrase. Le texte
s'arrête aux quatre cinquièmes de la page, au milieu d'une phrase ; la
suite est à la vue 454.}

\page{454}

point $\underline{i}$, la courbe nodale se trouve entre la cinquième ligne fictive
parallelle a $\underline{ab}$ et le bord dela plaque. il est clair qu'alors
la plus grande valeur de $f(y)$ est $>$ que celle de $f(x)$, desorte
que la courbe doit se diriger parallellement à l'axe
des $\underline{x}$ pour rentrer dans l'intervalle où se trouve le point
$\underline{e}$ de départ.

On voit que cette seconde courbe nodale est beaucoup
plus resserrée que la première, puisqu'elle \add{ne} passe \uncertain{que} par
4 points $\underline{i}$, tandis que la première passe \struck{par} par six des
mêmes points.

Je ne m'arrêterai pas plus longtems à la description de
la présente figure nodale, parceque la seconde partie
dela même figure est \struck{tou} entièrement semblable à la première;
\struck{et} c'est ce qui a toujours lieu \struck{quelque} \add{quelle que} soit la valeur de
l'angle $\underline{\omega}$; \struck{mais} \add{et même} si cette valeur donne lieu à un nombre
pair de lignes fictives, la figure nodale est composée
de quatre parties egales et semblables \add{entr'elles}, desorte qu'il suffit
alors de décrire le quart dela figure.

Il est clair que l'ensemble des quatre courbes rentrantes
et dela diagonale qui forme la figure nodale
complette, passe comme le veut la théorie par les
$\underline{n}^2 = 25$ points de repos invariables.

5.\textsuperscript{eme} valeur de l'angle $\underline{\omega}$. A partir du point $\underline{a}$ je passe
\note{En tête, au milieu, « (48) », non barré ; aucun nombre en haut à
droite : c'est le dos du feuillet de la vue 453. La première ligne
achève la phrase de la vue 453 (« Après avoir passé par un / point
$i$ »). « ne » est écrit au-dessus de « passe » ; le mot qui suit est
noirci, et « que » est offert avec doute. « par » est écrit deux fois
devant « six », le premier barré. « tou » est barré devant
« entièrement » ; le « et » qui ouvre la ligne suivante est barré ;
« quelque » est barré et « quelle que » écrit au-dessus ; « mais » est
barré et « et même » écrit au-dessus ; « entr'elles » est ajouté
au-dessus de « desorte ». « 5.\textsuperscript{eme} » est écrit en
exposant devant « valeur de l'angle $\omega$ », qui fait titre. Le texte
s'arrête aux trois quarts de la page, au milieu de la phrase, qui se
poursuit à la vue 455 (« je passe / comme à l'ordinaire audelà »).}

\page{455}

comme à l'ordinaire audelà dela seconde ligne fictive parallelle à
$\underline{ab}$ jusqu'au point $\underline{e}$ auquel appartient la plus grande valeur
de $f(y)$ dans cette intervalle\struck{s}; et parceque cette valeur est $<$ que
celle de $f(x)$, la portion dela courbe qui y est comprise, est
dirigée parallellement à l'axe des $y$. En sortant de cet intervalle,
la courbe traverse \struck{d'un côté le} celui où elle rencontre la
diagonale qu'elle coupe, par une ligne presque droite. \struck{Ensuite}
après avoir passé \add{par un point $i$, elle se trouve} en deça dela première ligne fictive; \struck{et} comme
alors la plus grande valeur de $f(y)$ est $>$ que celle de $f(x)$,
la courbe se dirige parallellement à l'axe des $\underline{x}$.
Elle reste ensuite dans l'intervalle entre la première et
la seconde lignes fictives; et elle se dirige parallellement à
l'axe des $y$, parceque la plus grande valeur de $f(y)$ est $<$
que celle de $f(x)$. Après avoir encore passé par un nouveau
point $i$, la courbe se trouve une seconde fois dans un même
intervalle avec la diagonale qu'elle coupe par une ligne
presque droite. En sortant de cet intervalle elle se trouve
entre la troisième et la quatrième lignes fictives, et elle se
dirige parallellement à l'axe des $y$ pour rentrer dans
l'intervalle où est situé le point de départ.

D'après la remarque que j'ai faite dans le cas précédent,
il est inutile de pousser plus loin la description de
la figure nodale; et il est clair qu'elle doit être formée
de quatre courbes rentrantes semblables à \struck{celle} la première.
On voit que l'ensemble de ces quatre courbes et des deux
diagonales nodales qui forme la figure nodale complette
\note{En tête, au milieu, « (49) », barré d'un trait horizontal ; en haut
à droite, à l'encre, « 223 ». La première ligne achève la phrase de la
vue 454 (« 5.\textsuperscript{eme} valeur de l'angle $\omega$. A partir du point
$a$ je passe / comme à l'ordinaire audelà »). « intervalle » est écrit « cette
intervalle », avec une lettre finale barrée. Les signes « $<$ » et
« $>$ » sont écrits tels quels dans la phrase (« est $<$ que celle de
$f(x)$ »). « d'un côté le » est barré devant « celui » ; « Ensuite » est
barré en fin de ligne ; « par un point $i$, elle se trouve » est ajouté
au-dessus de « en deça » ; « et » est barré devant « comme » ; « celle »
est barré devant « la première ». La page se termine sur « complette »,
sans point ; la phrase se poursuit à la vue 456 (« complette / passe
comme le veut la théorie »).}

\page{456}

passe comme le veut la théorie par les $\underline{n}^2 = 36$ points de
repos invariables.

Les cas suivans n'ayant pas été soumis à l'experience, je ne
m'y arreterai pas; mais il est aisé de voir que l'on parviendroit
à décrire très exactement les figures qui leur conviendroient,
si on prenoit la peine \struck{de calculer} et de comparer \add{entr'elles} les differentes
valeurs de $f(x)$ et $f(y)$ qui appartiennent, dans chaque cas,
aux points renfermés dans les differens intervalles des lignes
fictives, valeurs qui ne sont autre chose que celles de $\underline{z}$ pour
la lame dont les extrêmités sont libres au sens d'Euler.

Il est clair, aureste, que chacune de ces figures, excepté
\add{pourtant} celle qui se rapporte à la troisième valeur de l'angle $\omega$,
contiendra une courbe nodale rentrante\struck{, dont} \add{assujettie aux conditions suivantes}: le sommet \struck{de}
\struck{cette courbe} sera dans l'intervalle \struck{de la} \add{entre les} seconde et troisième lignes fictives;
\struck{que la portion dela courbe qui sera dans cet intervalle se dirigera} \add{et la courbe se dirigera dans cet intervalle}
parallellement à l'axe des $y$ parceque la plus grande valeur
de $f(y)$ est \add{alors} $<$ que celle de $f(x)$; \struck{que elle courbe} \add{la même courbe} \struck{traversera}
diagonalement l'intervalle dans lequel elle rencontrera la diagonale
parceque dans cet intervalle la plus grande valeur de $f(y) =$ celle de $f(x)$;
\struck{qu'ensuite} elle se prolongera dans une direction parallelle à l'axe
des $x$ parceque la plus grande valeur de $f(y)$ \struck{est} \add{sera} $>$ que celle de $f(x)$
\struck{qu'ensuite} dans l'intervalle entre la \struck{seconde et} première et la seconde
lignes fictives, \struck{elle} \add{la courbe dans cette ligne} se dirigera parallellement à l'axe des $y$ \struck{par les}
\struck{raisons pareilles}; \struck{et} \struck{qu}'elle traversera une seconde fois, dans une
direction diagonale, \add{un} \struck{l'}intervalle où elle rencontrera la diagonale
\note{Le feuillet, plus petit, est monté bas sur sa feuille. En tête, au
milieu, « (50) », non barré ; aucun nombre en haut à droite : c'est le
dos du feuillet de la vue 455. La première ligne achève la phrase de la
vue 455 (« qui forme la figure nodale complette / passe comme le veut
la théorie »). « de calculer » est barré ; « entr'elles » est ajouté
au-dessus de « comparer ». « pourtant » est écrit dans la marge de
gauche devant « celle qui ». « , dont » est barré et « assujettie aux
conditions suivantes » écrit au-dessus, d'une écriture fine ; le « L »
de « le sommet » est une capitale repassée. « de cette courbe » est
barré, en partie dans la marge ; « de la » barré et « entre les » écrit
au-dessus. La ligne « que la portion dela courbe qui sera dans cet
intervalle se dirigera » est barrée, et « et la courbe se dirigera
dans cet intervalle » est écrit au-dessus, le « elle » de « et elle »
barré et « courbe » écrit dessus. « alors » est ajouté au-dessus de
« $<$ » ; « que elle courbe » est barré et « la même courbe » écrit
au-dessus ; « traversera » est barré, sans que la phrase soit refaite.
« qu'ensuite » est barré deux fois ; « est » est barré et « sera » écrit
au-dessus ; « seconde et » est barré avant « première ». « elle » est
barré et « la courbe dans cette ligne » écrit au-dessus, lecture que le
sens ne rend pas sûre pour « ligne ». « par les raisons pareilles »,
« et » et « qu » sont barrés ; « l' » est barré devant « intervalle » et
« un » écrit au-dessus. La phrase se poursuit à la vue 457 (« la
diagonale / nodale; »).}

\page{457}

nodale; \struck{et qu}'elle changera encore une fois de direction avant de rentrer
\struck{dans les divers intervalles des lignes fictives valeurs qui restent autre chose}
dans l'intervalle où se trouve son sommet, \struck{desorte que}
\struck{que celles de $\underline{z}$ pour la lame dont les extrêmités sont libres \ill{}}

\struck{un seul \ill{}} Cette courbe passera par \uncertain{12} des points de repos
invariables. \struck{Dans le cas} \add{\uncertain{lorsqu'il s'agit}} dela 3.\textsuperscript{eme} valeur de $\omega$ (\struck{on} voir fig. 72 a) que la courbe

\marginal{rentrante passe par
8 des mêmes points $i$
parceque dans ce
cas seulement elle
\add{rencontre} \struck{traverse quatre fois}
les deux diagonales
nodales.}

N.o 12 Indépendamment des figures qui viennent d'être décrites,
il peut en exister plusieurs autres qui soient comprises dans
la même intégrale (C) et appartiennent aux mêmes valeurs
de l'angle $\underline{\omega}$; il suffit pour \struck{cela} \add{les expliquer} de changer le rapport
entre les constantes $\mathcal{A}$ et $B$ qui multiplient $f(x)$ et $f(y)$
dans la valeur de $z$. On voit en effet que, pour tout
autre rapport que celui exprimé par l'équation $\mathcal{A} = -B$,
il ne peut y avoir de ligne nodale diagonale, \struck{et} \add{et} que \struck{tout}
la figure \struck{que} et la situation des autres lignes nodales que
l'on a trouvées dans cette supposition, doivent changer en
même tems qu'elle\struck{, et} \add{On voit aussi} que les points invariables de repos \struck{sont} étant
les seuls qui soient communs à toutes les figures nodales qui
se rapportent a une même valeur de l'angle $\underline{\omega}$, \struck{et que} ces
figures \struck{pourront} \add{peuvent} être composées de lignes sinueuses auxquelles
les lignes fictives \struck{servoient} \add{serviront} d'axes, et qui couperont \add{ces mêmes} lignes
dans \struck{les} ces points invariables de repos.

Si on nomme avec M.r Chladni (V. N.o 109 de son traité d'acoustique)
\emph{flexion} « toute la déviation d'une ligne nodale vers un côté, consistant
« dans un éloignement et un rapprochement dela ligne droite laquelle
« on peut s'imaginer comme forme originaire » (cette droite est \add{véritablement} celle que je nomme
\marginal{mais M.r Chladni par la définition qu'il donne de la \emph{flexion} paroit \struck{la placer} \uncertain{tirer} par les points
les plus hauts ou les
plus bas dela ligne
sinueuse tandis qu'elle
doit passer par les
points moyens et servir d'axe a cette courbe)}
ligne fictive). on trouvera qu'il doit y avoir une \emph{flexion} pour
la première valeur de l'angle $\underline{\omega}$, une \emph{flexion} $+\frac{1}{2}$ pour la
\note{En tête, au milieu, « (51) », barré d'un trait horizontal ; en haut
à droite, à l'encre, « 224 ». Le premier mot, « nodale; », achève la
phrase de la vue 456 (« rencontrera la diagonale / nodale; »), dont la
suite est reprise et barrée ici. Les quatre premières lignes sont très reprises : « et
qu » est barré devant « 'elle changera » ; la deuxième et la quatrième
lignes sont barrées d'un trait horizontal (la fin de la quatrième, après
« libres », ne se lit pas), « desorte que » au bout de la troisième
aussi ; deux grands traits obliques en croix traversent en outre les
lignes deux à quatre, sans que « dans l'intervalle où se trouve son
sommet, » soit barré d'un trait propre. À la cinquième ligne, « un
seul » et un mot noirci qui le suit sont barrés ; après « par », deux
signes surchargés sont lus avec doute « 12 ». « Dans le cas » est barré,
et les mots écrits au-dessus sont lus avec doute « lorsqu'il s'agit » ;
la phrase ainsi corrigée n'est pas raccordée à ce qui précède ; la phrase se poursuit dans la marge de gauche (« rentrante
passe par 8 des mêmes points $i$ … nodales »), où « traverse quatre
fois » est barré et « rencontre » écrit au-dessus. Dans la parenthèse,
« on » est barré et « voir », sur un mot surchargé, n'est pas sûr. Le
« N.o 12 » est dans la marge, en tête de l'alinéa, à la suite du
« N.o 11 » de la vue 443. « les expliquer » est écrit au-dessus de
« cela », barré ; un second « et » au-dessus du premier, barré ;
« On voit aussi » au-dessus de « , et », barré ; « peuvent » au-dessus
de « pourront » ; « serviront » au-dessus d'un premier mot barré, lu
« servoient » ; « mêmes » au-dessus de « ces lignes ». La citation de
Chladni porte les guillemets en tête de chaque ligne, et « flexion »
est souligné ; « véritablement » est écrit au-dessus de « celle », sur
un mot surchargé, et sous « que je nomme » quelques mots très fins, peut-
être « la ligne sinueuse », ne sont pas rattachés au texte. La remarque
« mais M.r Chladni … a cette courbe) » est écrite dans l'interligne,
sous « que je nomme », puis continue dans la marge de gauche ; elle est
donnée ici comme note marginale, à l'endroit où elle commence ; « la
placer » y est barré et « tirer » est offert avec doute. La page se
termine sur « pour la » ; la vue 458 continue.}

\page{458}

seconde, 2 \emph{flexions} pour la troisième, 2 \emph{flexions} $+\frac{1}{2}$ pour la
quatrième et enfin 3 \emph{flexions} pour la 5.\textsuperscript{eme} valeur du même
angle $\underline{\omega}$; car dans le premier \add{cas} il y a 2 points de repos
invariables sur la ligne fictive, dans le second il existe 3
points semblables, dans le troisième il y en a 4, dans
le quatrième, \struck{cinq} enfin dans le cinquième cas il
y a six \add{de ces} points \struck{de pareil nature} et l'addition d'un point
de repos invariable donne une demi flexion, parceque la sinuosité terminée
\marginal{par deux points consécutifs
ne forme, suivant M.r Chladni
qu'une demi flexion excepté
pour la 1.\textsuperscript{ere} valeur de
l'angle $\underline{\omega}$.}
Il est aisé de voir maintenant combien \struck{cette} la théorie \struck{est} \uncertain{que j'ai exposée}
est conforme à l'expérience. En effet la fig 64 de M.r
Chladni est précisément celle que j'ai trouvée pour la
première valeur de l'angle $\underline{\omega}$ dans la supposition
$\mathcal{A} = -B$. La fig. 65. \struck{montre} qui appartient à la même
valeur de l'angle $\omega$ \add{(pour \struck{cette} ce cas et les suiv. voyez la table p. 150 de Chladni, dans laquelle
il indique les sons qui répondent a chaque fig.)} montre une seule \emph{flexion} avec cette
particularité que les deux lignes infléchies auxquelles les
deux lignes fictives servent d'axe, se rejoignent très près
du bord dela plaque pour former une courbe rentrante.
Il est clair que c'est du rapport entre $\mathcal{A}$ et $B$ que
depend\struck{ent} l'élévation de chaque sinuosité au dessus des
lignes fictives, desorte que cette ligne même est une limite
tandis que la courbe rentrante est l'autre limite de
l'écartement de la ligne nodale de sa situation
primitive.
\note{En tête, au milieu, « (52) », non barré ; aucun nombre en haut à
droite : c'est le dos du feuillet de la vue 457. La première ligne
achève la phrase de la vue 457 (« une \emph{flexion} $+\frac{1}{2}$ pour
la / seconde »). « flexion(s) » est souligné chaque fois. « cas » est
ajouté au-dessus de « premier » ; « cinq » est barré devant « enfin » ;
« de ces » est écrit au-dessus de « six », et « de pareil nature » barré.
La ligne « de repos invariable donne une demi flexion, parceque la
sinuosité terminée » est serrée entre deux lignes, d'une écriture plus
petite, et la phrase continue dans la marge de gauche (« par deux points
consécutifs … l'angle $\omega$. ») ; la note est placée ici à l'endroit
où la phrase y passe. Devant « théorie », « cette » est barré et « la »
écrit dessus ; après, « est » est barré et « que j'ai exposée », d'une
écriture très fine au bout de la ligne, n'est pas assuré. « montre » est
barré après « fig. 65. » ; le renvoi à la table de Chladni est écrit
au-dessus de la ligne suivante, sur deux lignes très fines, et « cette »
y est barré devant « cas ». Le $\mathcal{A}$ de « entre $\mathcal{A}$ et
$B$ » est surchargé. La fin de « dependent » est barrée. Le texte
s'arrête au tiers inférieur de la page ; le reste est blanc.}

\page{459}

La figure 67. appartient à la seconde valeur de l'angle $\underline{\omega}$, dans
la supposition $B=0$. La fig. 67 c est très exactement celle que
j'ai décrite pour la même valeur de $\underline{\omega}$ dans la supposition
$\mathcal{A} = -B$; et la fig. 67 b qui appartient encore à la même valeur
de l'angle $\underline{\omega}$, montre une \emph{flexion} $+\frac{1}{2}$ autour de la position
de chacune des lignes qui sont droites dans la figure primitive
fig. 67.

La fig. 72 a est celle que j'ai décrite pour la troisième valeur de
l'angle $\underline{\omega}$, la fig. 72 b qui appartient à la même valeur de l'angle
$\underline{\omega}$ montre 2 \emph{flexions} autour de chacune des lignes fictives, c'est-à-dire
autour des lignes qui seroient droites dans la supposition $B=0$
la fig. 73 a qui appartient aussi à la même valeur de l'angle $\underline{\omega}$
montre egalement 2 \emph{flexions}: mais il arrive ici, comme dans la
fig. 65, que les lignes sinueuses se rejoignent. \struck{(je ne sais pourquoi M.r Chladni a changé
le chiffre de cette fig. et de la suivante
\uncertain{ce qui} peut \uncertain{faire croire} que les mêmes \ill{})} Enfin la fig 73 b que
appartient encore a la 3.\textsuperscript{eme} valeur de l'angle $\underline{\omega}$, montre
aussi 2 \emph{flexions}. On voit qu'il peut exister un très grand
nombre de figures differentes pour chaque valeur de l'angle $\underline{\omega}$:
mais l'experience n'en a montré jusqu'ici qu'une petite partie; et
\struck{souvent} même les plus simples manquent dans la plus part des
cas, \struck{c'est ainsi que la fig. 78 qui appartient} \struck{et celle que j'ai trouvé} \struck{à la quatrième}
\struck{par ex} la figure primitive, par exemple, ne se realise que dans
\struck{valeur de l'angle $\underline{\omega}$, est la seule que l'experience donne}
un seul \add{de ces} cas, \struck{fig. 6} savoir pour la 2.\textsuperscript{eme} valeur de l'angle $\omega$
\struck{pour ce cas, desorte que les lignes nodales droites dependantes}
(V. fig 67 de Chladni)
\struck{de la seule supposition $B=0$. ne se montrent ni pour le cas} \struck{dont il s'agit} \struck{ni}
\struck{pour la plus part des autres}.
\note{En tête, au milieu, « (53) », barré d'un trait horizontal ; en haut
à droite, à l'encre, « 225 ». La page s'ouvre sur un alinéa nouveau ; la
vue 458 s'achevait sur « primitive. ». Les numéros de figures sont ceux
de l'ouvrage de Chladni, comme annoncé à la vue 445 (« je désignerai
chaque figure par le chiffre que porte son analogue dans l'ouvrage de
M.r Chladni ») : 64, 65, 67, 67 b, 67 c, 72 a, 72 b, 73 a, 73 b. Dans
« fig. 72 b », le 2 est surchargé sur un autre chiffre. La note entre
parenthèses, écrite en trois lignes très fines au bout de la ligne de
« se rejoignent », est barrée ligne à ligne ; sa troisième ligne est lue
avec doute et sa fin ne se lit pas. Dans la marge de droite, en face,
quelques lettres verticales, peut-être « V. fig. », ne sont pas lues. Le
bas de la page est une suite de rédactions barrées, enchevêtrées ; ce
qui n'est pas barré se lit : « et même les plus simples manquent dans la
plus part des cas, la figure primitive, par exemple, ne se realise que
dans un seul de ces cas, savoir pour la 2.\textsuperscript{eme} valeur de
l'angle $\omega$ (V. fig 67 de Chladni). » L'ordre des passages barrés
donné ici est celui des lignes ; « fig. 78 » et « dependantes » ne sont
pas assurés ; « et celle que j'ai trouvé » (au-dessus de « appartient »)
et « dont il s'agit » (au-dessus de « le cas ») sont des additions
elles-mêmes barrées. Un point final est posé après la dernière ligne
barrée. La page s'arrête aux quatre cinquièmes ; la vue 460 ouvre un
alinéa nouveau.}

\page{460}

A l'égard des figures 78 et 85 qui se rapportent aux 4.\textsuperscript{eme}
et 5.\textsuperscript{eme} valeurs de l'angle $\underline{\omega}$, On voit que j'ai marqué une
partie des lignes nodales que j'ai trouvé\struck{es} devoir composer les
figures nodales, par des lignes pleines tandis que les autres
sont indiquées par des lignes ponctuées \struck{et qu'une} portion de
ces lignes ponctuées \struck{est traversée} par des lignes pleines qui
rejoignent les lignes pleines des \struck{la} figures que j'ai décrites
et \add{qui} composent avec elles des figures nodales semblables
a celles que M.r Chladni a trouvées et qu'il a numérotées
par des mêmes chiffres.

Au premier coup d'œil, on seroit tenté de croire que, dans les
cas dont il s'agit ici, la théorie et l'experience ne sont pas d'accord:
mais un instant de reflexion suffit pour les concilier. En effet,
si on adopte \struck{la théorie que j'ai} se rapelle ceque j'ai cherché
a établir \uncertain{§ 3}, on verra que toute\struck{s} ligne qui satisfait aux
conditions des limites, pouvant être considérée comme une veritable
extrêmité analytique, les parties qu'elles séparent peuvent vibrer
comme si elles étoient reellement isolées; pourvu toutefois que
le son de l'une soit le même que celui de l'autre; car
sans cette condition, le mouvement qui dans la plaque entière
se communique d'une partie a l'autre, deviendroit impossible.
Or, si les portions dela plaque\struck{s} étoient isolées, elles rendroient
toujours le même son lorsque leurs mouvemens se
\note{Le feuillet, plus petit, est monté bas sur sa feuille. En tête, au
milieu, « (54) », non barré ; aucun nombre en haut à droite : c'est le
dos du feuillet de la vue 459, numéroté 225. La page s'ouvre sur un
alinéa nouveau, la vue 459 s'achevant sur un point. Les figures 78 et
85 sont celles de Chladni. Dans « trouvées », la fin est barrée ; « et
qu'une » est barré devant « portion », et « est traversée » après « ces
lignes ponctuées », de sorte que la phrase se lit sans verbe, telle
qu'elle est laissée. « la » est barré après « des » ; « qui » est
ajouté au-dessus de « composent ». Un cadre fin, tracé au crayon, entoure
une partie du premier alinéa. « la théorie que j'ai » paraît barré d'un
trait long ; la lecture « si on adopte se rapelle » est celle de la
page, sans le « et » ou le « qu'on » que le sens attend. Devant « 3 »,
un signe haut et étroit est lu comme le § de son écriture ; « § 3 »
renverrait au paragraphe des « propriétés analytiques des differens
points », mais ce n'est pas assuré. L'« s » de « toutes » est barré, de
même que celui de « plaques ». « isolées » s'accorde avec « parties ».
La page s'arrête aux quatre cinquièmes, au milieu de la phrase, sur
« mouvemens se » ; la vue 461, hors de ce lot, commence par
« rapporteroient à la même valeur de l'angle $\omega$ ».}

\end{document}
